% docs/manual/operations/chapters/02-hpc-install.tex
% 第 2 章：HPC 安装与依赖管理

\chapter{HPC 安装与依赖管理}

HPC 集群的 Python 环境是 \openbench{} 部署中最折磨人的一环：受限网络、共享只读、module 系统、login vs 计算节点的 Python 不一致。本章给四条可行路径，按你环境的限制选一条。

\section{HPC 环境变体}

主流 HPC 集群提供以下几种 Python：

\begin{itemize}
  \item \textbf{module load python}：管理员预编译；版本固定；\texttt{python --version} 是这个
  \item \textbf{共享 conda}：管理员维护一个全局 conda env，用户激活即可
  \item \textbf{容器（Singularity / Apptainer）}：自带完整运行时；\openbench{} 装在容器内
  \item \textbf{用户自管}：自己装 conda / uv / pyenv 在 \texttt{\$HOME}
\end{itemize}

绝大多数集群至少有第一种与第四种。本章四节对应四种路径。

\section{路径 A：module load + 用户级 venv}

最常见的"轻量"方案。管理员维护 \texttt{python/3.12.x}，你自己建 venv 装 \openbench{}。

\begin{minted}{bash}
# 1. 加载 module（具体名字看你的集群文档）
module load python/3.12.0

# 2. 建 venv
python -m venv ~/openbench-venv
source ~/openbench-venv/bin/activate

# 3. pip 装
pip install --upgrade pip
pip install "openbench[remote,report]"   # 注意不装 [gui]

# 4. 验证
openbench version
\end{minted}

\warninline{HPC 上 pip 可能需要 \texttt{--user} 或 venv，因为 module 加载的 site-packages 通常只读。优先 venv。}

\section{路径 B：用户自管 conda}

如果系统提供的 Python 版本太旧（< 3.10），用 miniforge：

\begin{minted}{bash}
# 1. 装 miniforge（社区维护，纯 conda-forge）
wget https://github.com/conda-forge/miniforge/releases/latest/download/Miniforge3-Linux-x86_64.sh
bash Miniforge3-Linux-x86_64.sh -b -p $HOME/miniforge3
source $HOME/miniforge3/bin/activate

# 2. 建 env
conda create -n openbench python=3.12 -y
conda activate openbench

# 3. cartopy 系统依赖通过 conda-forge 解决（绕过 pip 编译）
conda install -c conda-forge cartopy

# 4. pip 装其余
pip install "openbench[remote,report]"
\end{minted}

优点：依赖最稳。缺点：占空间多（~3-5 GB）。

\section{路径 C：受限网络的离线安装}

某些 HPC 出口网络受限，pip install 直接失败。流程：

\subsection{在有网的机器上}

\begin{minted}{bash}
# 把 openbench 与全部依赖打包成 wheelhouse
mkdir openbench-wheelhouse
pip download "openbench[remote,report]" -d openbench-wheelhouse/

# 也下载 openbench 本身（如要装本地源码）
pip download . --no-deps -d openbench-wheelhouse/

# 打包传输
tar czf openbench-wheelhouse.tar.gz openbench-wheelhouse/
\end{minted}

\subsection{在 HPC 上}

\begin{minted}{bash}
# 1. 解包
scp openbench-wheelhouse.tar.gz hpc:~/
ssh hpc
tar xzf openbench-wheelhouse.tar.gz

# 2. 离线安装（不联网）
module load python/3.12.0
python -m venv ~/openbench-venv && source ~/openbench-venv/bin/activate
pip install --no-index --find-links=./openbench-wheelhouse "openbench[remote,report]"
\end{minted}

\textbf{注意点}：

\begin{itemize}
  \item wheelhouse 必须用 \emph{相同 Python 版本} 与 \emph{相同平台}（manylinux2014\_x86\_64 或类似）下载
  \item cartopy 不是 pure Python，需匹配的 wheel；conda-forge 有最全的预编译
  \item 第一次成功安装后，可把 \texttt{\~{}/openbench-venv/} 整体打包给同事复用
\end{itemize}

\section{路径 D：容器（Singularity / Apptainer）}

最隔离的方案：build 一个含 \openbench{} 与依赖的 container image，集群上用 singularity 运行。

\begin{minted}{bash}
# 在有 sudo 的机器上 build（HPC login 节点通常无 sudo）
cat > openbench.def <<EOF
Bootstrap: docker
From: python:3.12-slim

%post
    apt-get update && apt-get install -y libgeos-dev libproj-dev proj-bin
    pip install "openbench[remote,report]"

%runscript
    exec openbench "$@"
EOF

singularity build openbench.sif openbench.def

# scp openbench.sif 到 HPC
# 之后 HPC 上：
singularity exec openbench.sif openbench run config.yaml
\end{minted}

\textbf{优点}：跨集群复用、版本完全可复现。
\textbf{缺点}：build 阶段需要 root；容器镜像 1-2 GB。

\section{共享只读环境 + 用户私有 cache}

许多 HPC 倾向于"集群管理员维护一个公共 \openbench{} 安装；用户共享，但 cache 与输出走自己 \$HOME"。

\subsection{管理员侧}

\begin{minted}{bash}
# /opt/openbench/3.0.0/  （只读）
sudo mkdir -p /opt/openbench/3.0.0
cd /opt/openbench/3.0.0
python -m venv venv
source venv/bin/activate
pip install "openbench[remote,report]"

# 加 module file 暴露给用户
\end{minted}

\subsection{用户侧}

\begin{minted}{bash}
module load openbench/3.0.0
openbench version

# 用户的 yaml 与输出走自己 HOME
cd ~
cat > openbench.yaml <<EOF
project:
  name: my-eval
  output_dir: ~/openbench-output    # 写入用户目录
  ...
EOF
openbench run openbench.yaml
\end{minted}

\openbench{} 的 platformdirs 会把 user-level catalog 放到 \texttt{\$HOME/.config/openbench/} 或对应位置；不污染管理员的全局安装。

\warninline{Registry 的 \emph{user-level catalog} 默认位置因系统不同：macOS \texttt{\~{}/Library/Application\ Support/openbench/}，Linux \texttt{\~{}/.config/openbench/}。HPC 上如果 \$HOME 配额很紧，可设环境变量 \texttt{XDG\_CONFIG\_HOME} 重定向。}

\section{Login 节点 vs 计算节点}

\textbf{陷阱}：很多 HPC 的 login 节点与计算节点 Python 版本/路径不一致。在 login 装好 \openbench{}，提交到计算节点时找不到。

\subsection{诊断}

\begin{minted}{bash}
# 在 login 节点
which python
python --version

# 提交一个交互节点跑同样命令
srun --pty bash
which python
python --version
\end{minted}

如果两者不同，必须改用：

\begin{itemize}
  \item Conda：把 conda activate 写到 SLURM 脚本里
  \item 容器：用 singularity exec 一致跑
  \item 自管 venv：用绝对路径 \texttt{/path/to/venv/bin/openbench}
\end{itemize}

\subsection{SLURM 脚本模板}

\begin{minted}{bash}
#!/bin/bash
#SBATCH --job-name=openbench-eval
#SBATCH --cpus-per-task=32
#SBATCH --mem=128G
#SBATCH --time=24:00:00

# 关键：在 SLURM job 内重新激活环境
source ~/miniforge3/bin/activate openbench

# 与 SLURM 分配核数一致
export OPENBENCH_NUM_CORES=$SLURM_CPUS_PER_TASK

cd /scratch/$USER/openbench-runs
openbench run config.yaml --cores $SLURM_CPUS_PER_TASK
\end{minted}

\section{不要装 [gui] 的原因}

GUI extra 会拉 \modname{PySide6} —— 体积大（~200 MB）+ 依赖系统级 Qt 库，HPC 上常装不上：

\begin{minted}{text}
ERROR: Failed building wheel for PySide6-Essentials
ERROR: Could not build wheels for PySide6-Essentials
\end{minted}

即使强行装上，HPC 计算节点没有图形系统，启动 GUI 会报：

\begin{minted}{text}
qt.qpa.plugin: Could not load the Qt platform plugin "xcb"
\end{minted}

GUI 是本地工具：在你笔记本写好 yaml，scp 到 HPC，用 CLI 跑。

\section{升级 OpenBench}

\begin{minted}{bash}
source ~/openbench-venv/bin/activate
pip install -U openbench

# 升级后第一次 run 可能要清旧 cache
rm -rf ~/openbench-output/<old-run>/scratch
\end{minted}

升级前后建议跑一次 \cli{openbench check} 确认配置仍合法（schema 变化时会有 deprecation 警告）。

\section{验证安装}

每次新安装后跑这套：

\begin{minted}{bash}
openbench version              # 版本号正常
openbench --help               # 7 个子命令都列出来
openbench data list            # registry 加载成功
openbench model list           # model profile 21 个

# 用一个最小配置跑 check
cat > /tmp/test.yaml <<EOF
project:
  name: smoke
  output_dir: /tmp/openbench-smoke
  years: [2010, 2010]

evaluation:
  variables: [Evapotranspiration]

reference:
  Evapotranspiration: GLEAM_v4.2a

simulation:
  test:
    model: CoLM2024
    root_dir: /nonexistent
EOF

openbench check /tmp/test.yaml
# 预期：YAML 与 schema 验证通过；reference 解析成功；
# 不会真去查 simulation 数据（root_dir 路径不存在不算 check 错误）
\end{minted}

如果上述命令任一失败，根据用户卷第 1 章 / 第 10 章排查。

\section{下一步}

下一章详解远程 SSH 执行：从笔记本提交评估到 HPC 节点的完整流程。
